\documentclass[10pt,a4paper]{article}
\usepackage[T1]{fontenc}
\usepackage{lmodern,amsmath,amssymb,amsthm,graphicx,booktabs,geometry,hyperref}
\geometry{margin=2.3cm}
\hypersetup{colorlinks=true,urlcolor=blue,linkcolor=blue,citecolor=blue}
\setlength{\emergencystretch}{2em}
\newtheorem{proposition}{Proposition}
\newtheorem{lemma}{Lemma}
\newcommand{\V}{V_h}
\newcommand{\W}{W_k}
\newcommand{\PV}{P_V}
\newcommand{\PW}{P_W}
\title{Energy-compatible memory coupling for semi-Lagrangian finite elements\\on changing auxiliary spaces}
\author{Rodin / SLFE research draft}
\date{17 November 2026}
\begin{document}
\maketitle
\begin{abstract}
A fixed finite element observation is coupled to a persistent field on a
moving auxiliary transport space. Support changes preserve the finite
element observation while compressing the unresolved component by an
orthogonal projection. The projection space may depend on the state:
pairwise non-expansivity can fail, but an exact relative-energy identity
still controls distance to every finite element reference. This yields an
error estimate containing local evolution defects without an additional
term proportional to the number of support changes. Complete-state
reaction moments provide a dissipative physics update; reacting only the
visible finite element component can leave a substantial hidden residue.
Shared compression spaces preserve linear relations between transported
components. GPU experiments examine strong reaction, frequent remapping,
conditioning and complete computational cost. A rational orthogonal
embedding is also tested as a linear alternative with a positive energy
representation. The analysis is conditional on the transport defects;
a general PDE convergence rate, general finite element velocity fields
and bibliographic novelty remain open.
\end{abstract}
\noindent\textbf{Status.} Working manuscript, not submitted. The experiments
are isolated CUDA prototypes; this is not a production integration into
Rodin's variational language. Full parameter records, unsuccessful avenues
and implementation controls accompany the illustrated French report.

\section{Scope and relation to existing methods}
Semi-Lagrangian finite element enrichment, dynamic subscales, conservative
remapping and particle--grid moment preservation are established subjects
\cite{ouardghi,codina,simon,adams,keith}. A fixed finite element mesh or the
storage of a background field plus finite element increments is therefore
not, by itself, a distinctive contribution. In particular, Remark 2.0.1 of
\cite{keith} stores a quadrature background and finite element increments.
The question considered here is narrower: how can an unresolved physical
memory be transported between auxiliary supports, remain compatible with
finite element physics, and admit useful stability and error statements
when its compression space depends on the solution?

The setting is a periodic domain, a fixed finite-dimensional subspace
$\V\subset L^2(\Omega)$ and a sequence of auxiliary spaces $W_n$.
All projections in the analysis are orthogonal in the same physical
$L^2$ inner product. The computational realization uses discontinuous
$P_p$ functions on affine triangles and tensor-product $Q_q$ auxiliary
functions. The present transport tests use prescribed velocities with
exact translations or compositions of separable shears. They do not
cover a general velocity supplied by a coupled finite element solve.

\section{A physical memory and its finite element observation}
Write $R=\PW|_{\V}$, so $R^*=\PV|_W$. The physical state and its
observable are
\begin{equation}
 q=a+d,\qquad a\in\V,\quad d\in W,\qquad v=\PV q=a+R^*d.
\end{equation}
It is convenient to store $a$ and $w=Ra+d=\PW q$. Then
\begin{equation}\label{eq:energy}
 q=a+w-Ra,\qquad
 \|q\|^2=\|w\|^2+\|a\|^2-\|Ra\|^2.
\end{equation}
Given a contraction $S:W\to W$, auxiliary advection acts as
\begin{equation}\label{eq:transport}
 \Phi_W q=S\PW q+(I-\PW)q,\qquad
 w^+=Sw,\quad v^+=v+R^*(Sw-w).
\end{equation}
The two components in the first expression are orthogonal, hence
$\|\Phi_W\|\leq1$. If $1\in W$ and $S$ preserves its integral,
the full state and its finite element observation preserve mass.
This argument identifies a contractive discrete evolution; consistency
with the continuous advection equation remains a separate requirement.

Likewise, a finite element contraction $G:\V\to\V$ has the extension
$G\PV+(I-\PV)$. This extension leaves the hidden component unchanged.
It is energy stable but is generally not the physical extension of a
reaction or diffusion operator to the unresolved field.

\section{Observation-preserving compression on a new support}
Let $W_+$ be the new auxiliary space and $H_+=(I-\PV)|_{W_+}$.
Choose directions $X=(x_1,\ldots,x_r)$ in $W_+$, possibly from the
current state. For a target $f\in L^2(\Omega)$, define
\begin{equation}\label{eq:compression}
 v_f=\PV f,\quad h_f=P_{W_+}f-R_+v_f,\quad
 G_X=X^*(I-R_+R_+^*)X,\quad b_X=X^*h_f.
\end{equation}
The candidate is
\begin{equation}\label{eq:project}
 \mathcal P_X f=v_f+H_+XG_X^\dagger b_X.
\end{equation}
The inverse acts only on a small Gram matrix. Formula
\eqref{eq:project} is the orthogonal projection onto
$\V\oplus\operatorname{range}(H_+X)$, including when the columns
are dependent. A spectral truncation instead projects onto the retained
hidden subspace; identity on discarded directions must not be assumed.

\begin{proposition}
For each realized choice of $X$,
\[
 \PV\mathcal P_X f=\PV f,\qquad
 \|f\|^2=\|\mathcal P_Xf\|^2+\|f-\mathcal P_Xf\|^2.
\]
If the target belongs to the retained space, the compression is exact.
\end{proposition}
\begin{proof}
The hidden range is orthogonal to $\V$. The Gram system is the normal
equation of the orthogonal projection onto that range. The claims follow
from orthogonality and Pythagoras.
\end{proof}

For support changes, take $f=q=a+d$ and $C=P_{W_+}|_W$. The moments
needed in \eqref{eq:compression} are available without reconstructing
a global function:
\[
 v_f=v,\qquad h_f=Cd+R_+(a-v).
\]
The experiments use the incoming direction $Cd$ and a few Krylov
directions generated from $h_f$ by $R_+R_+^*$. Augmented Krylov
constructions themselves are classical \cite{gaul}. At an unchanged
support, retaining the incoming direction retains the original state.

For several components, one common spatial space is formed from the
directions of every component. Since $\mathcal P_X$ fixes $\V$, a
relation $\sum_j c_j f_j=b\in\V$ implies
$\sum_j c_j\mathcal P_Xf_j=b$. The common space is unchanged by an
orthogonal rotation of the component basis when all generating
directions are rotated together. This gives component covariance of
the exact projection. Separate state-dependent spaces do not have
this property. Neither statement implies pointwise positivity.

\section{Relative energy and an error bound with changing supports}
The following identity requires less than pairwise non-expansivity.
It also applies to energy-reducing transfers that are not orthogonal
projections.

\begin{lemma}\label{lem:relative}
Let a possibly nonlinear map $B$ satisfy
$\PV B(q)=\PV q$ and $\|B(q)\|\leq\|q\|$. For every $u\in\V$,
\begin{equation}\label{eq:relative}
 \|B(q)-u\|^2=\|q-u\|^2-
 \bigl(\|q\|^2-\|B(q)\|^2\bigr).
\end{equation}
In particular, $\|B(q)-u\|\leq\|q-u\|$.
\end{lemma}
\begin{proof}
Expand both squared norms. The cross terms agree because
$\langle B(q),u\rangle=\langle\PV B(q),u\rangle
=\langle\PV q,u\rangle=\langle q,u\rangle$.
\end{proof}

\begin{proposition}\label{prop:error}
Let $q_{n+1}=B_n\Phi_nq_n$, with $B_n$ satisfying
Lemma~\ref{lem:relative} and $\Phi_n$ Lipschitz with constant
$\ell_n\leq1$. For arbitrary references $u_n\in\V$, set
$\delta_n=\|\Phi_nu_n-u_{n+1}\|$. Then
\begin{equation}\label{eq:bound}
 \|q_N-u_N\|\leq
 \left(\prod_{j=0}^{N-1}\ell_j\right)\|q_0-u_0\|
 +\sum_{n=0}^{N-1}
 \left(\prod_{j=n+1}^{N-1}\ell_j\right)\delta_n.
\end{equation}
The same bound holds for $\|\PV q_N-u_N\|$.
\end{proposition}
\begin{proof}
Apply the lemma with reference $u_{n+1}$, insert $\Phi_nu_n$ and
use the Lipschitz estimate. Iterating
$e_{n+1}\leq\ell_ne_n+\delta_n$ proves the result.
\end{proof}

The supports and compression spaces may be selected from the computed
trajectory: the argument uses each realized operator and reference.
There is no additional term proportional to the number of support
changes. This does not make remapping physically harmless. The chosen
supports affect every $\delta_n$, and the augmented norm can make the
bound loose. Pairwise distances between arbitrary hidden states can
increase under nonlinear compression without contradicting
Lemma~\ref{lem:relative}.

For an FE reference $\dot u=A_Vu$ and auxiliary generator $L_W$,
the infinitesimal defect of \eqref{eq:transport} has squared norm
\begin{equation}\label{eq:generator}
 \eta_W(u)^2=\|L_WRu-RA_Vu\|^2
             +\|A_Vu\|^2-\|RA_Vu\|^2.
\end{equation}
Both a commutator and an unresolved-generator term appear. Controlling
only the second term is insufficient. A PDE convergence theorem would
require estimates for these terms, temporal defects, and FE projection
error; it is not supplied by energy stability alone.

\section{Physical reaction and the failure of additive compression}
For $f=gq$, $g(x)=\exp(-\tau r(x))$, the required moments follow from
weighted local and cross-mass operators:
\begin{equation}
 \PV(gq)=F_ga+R_g^*d,\qquad
 \PW(gq)=R_ga+A_gd,
\end{equation}
where $F_g=\PV g|_{\V}$, $R_g=\PW g|_{\V}$ and
$A_g=\PW g|_W$. They define \eqref{eq:compression} for the physical
target, even though $gq$ need not belong to $\V+W$.
If $0\leq g\leq\gamma\leq1$, then
\[
 \|\mathcal P_X(gq)\|\leq\gamma\|q\|.
\]
The factor $\gamma$ also enters \eqref{eq:bound} through the
Lipschitz constant of multiplication by $g$. For a spatial scalar
rate times a constant symmetric component matrix, GPU diagonalization
of that matrix reduces the moment calculation to scalar modes. A
shared final spatial projection preserves conserved component sums.
Spatially varying eigenvectors and nonlinear reaction networks are
outside the present implementation.

An algebraically tempting alternative computes an increment
$\Delta q=(g-I)q$ and adds its compressed value. Let $P$ be the
realized orthogonal projector and $e=(I-P)q$. Then
\begin{equation}\label{eq:increment}
 q+P\Delta q=P(gq)+e,\qquad
 \|q+P\Delta q\|^2=\|P(gq)\|^2+\|e\|^2.
\end{equation}
Thus additive and full-target compression agree only if the old state
is retained. A rank cutoff can violate this requirement even when
the old hidden direction was nominally included. Strong damping makes
an absolutely small discarded component large relative to the final
physical solution. Reprojecting the complete target removes this
particular persistent residue.

\section{A linear alternative with a positive energy representation}
To avoid the subtractive energy in \eqref{eq:energy}, consider
$F=(R,I)^T$, $M=I+R^*R$ and the orthogonal graph reflection
\[
 U_R=2FM^{-1}F^*-I,\qquad
 \widehat J_Rv=U_R(0,v)^T
  =\bigl(2RM^{-1}v,(2M^{-1}-I)v\bigr)^T.
\]
An augmented state $(w,z)$ has energy $\|w\|^2+\|z\|^2$ and
observation $\widehat J_R^*(w,z)$. For a contraction
$C:W_0\to W_1$, the transition
\[
 \widehat B=U_{R_1}\operatorname{diag}(C,I)U_{R_0}
\]
is linear and contractive, satisfies
$\widehat J_{R_1}^*\widehat B=\widehat J_{R_0}^*$ and is the identity
when $R_1=R_0$, $C=I$. For a contraction $R$, every inverse has
spectrum in $[1,2]$. Fixed-degree Chebyshev evaluation avoids a
data-dependent stopping rule.

This is a different discretization: the first injection block is
$2R(I+R^*R)^{-1}$ rather than $R$. The physical reaction construction
above cannot simply be reused. The graph reflection is a classical
rational orthogonal parametrization, closely related to Cayley
coordinates \cite{cayley}; no linear-algebra novelty is claimed.

\section{GPU evidence and limitations}
All numerical assembly, solves, trajectory references and reductions
use double precision CUDA/CuPy on an RTX 5090. The CPU orchestrates
runs and generates plots. Dense GPU matrix oracles independently
check mass adjoints, observation preservation and projection identities
at degrees 1, 2 and 4. Compute Sanitizer logs accompany the data.
Reported speed comparisons concern the measured prototypes only.

\paragraph{Frequent support changes.}
Four sine--cosine fields of frequencies $(1,0)$ and $(1,2)$ are
translated with velocity $(1,0.37)$ to time $0.2$. Two perturbed
auxiliary partitions alternate. Rotating the initial FE coefficient
columns gives the exact FE reference, and the local defect norm is
independent of the Fourier phase. The 78 recorded configurations
respect \eqref{eq:bound}; the largest error-to-bound ratio is below
$0.20$. For 255 changes, the last measured spatial orders of the graph
variant are $1.93$, $2.99$ and $4.98$ for $P_1$, $P_2$ and $P_4$.
At $n=32$, the $P_4$ maximum component error changes from
$2.63\,10^{-7}$ at 255 changes to $2.95\,10^{-7}$ at 4095 changes.
These are observations for this smooth translation, not a general
order theorem or a claim that arbitrarily frequent remapping is free.
\begin{figure}[htbp]\centering
\includegraphics[width=\linewidth]{night-remap-consistency.pdf}
\caption{Translation with changing auxiliary supports. Physical errors
and the independently evaluated defect bound use distinct norms, as
specified in the illustrated report.}\label{fig:consistency}
\end{figure}

\paragraph{Strong spatial reaction.}
A transported localized profile reacts at rate
$\kappa(1+0.5\sin(2\pi x))$. Independent GPU characteristics integrate
both trajectory and reaction exponent. At $T=0.25$, $\kappa=100$,
the reference norm is approximately $4.20\,10^{-14}$; relative errors
are therefore essential. Table~\ref{tab:reaction} uses 256 time steps,
$P_4$ FE functions and auxiliary resolution $2n$.
\begin{table}[htbp]\centering
\begin{tabular}{llrr}\toprule
$n$ & Physics extension & Relative error & FE projection floor\\\midrule
64 & FE-only reaction & $6.05\,10^6$ & $2.37\,10^{-3}$\\
64 & Uniform hidden damping & $3.27\,10^1$ & $2.37\,10^{-3}$\\
64 & Complete physical moments & $2.40\,10^{-3}$ & $2.37\,10^{-3}$\\
128 & Complete physical moments & $1.39\,10^{-4}$ & $1.37\,10^{-4}$\\\bottomrule
\end{tabular}
\caption{Correcting the physical extension of reaction substantially
improves accuracy at additional computational cost. Absolute errors of
the incomplete models are small but exceed the strongly damped reference.}
\label{tab:reaction}
\end{table}
On the four-component exchange test at $n=64$, the complete moments give
relative error $7.80\,10^{-4}$ against an FE floor $7.80\,10^{-4}$,
compared with $7.70\,10^{-2}$ for FE-only reaction. The tested component
matrix is constant and symmetric. Shared spatial compression retains
the component sum to roundoff; no general positivity result is inferred.

\paragraph{Conditioning and genuine hidden residue.}
Positive quadrature of $|a+d|^2$ on the common FE--auxiliary intersections
avoids cancellation in \eqref{eq:energy}. Its degree is sufficient for
the square of the difference of $P_p$ and $Q_q$ polynomials. For two
constants separated by $10^{-12}$, it resolves an energy of order
$10^{-24}$ where subtractive formulas return roundoff-scale values.
In the strong-reaction additive experiment, the positive final energy
is $9.00\,10^{-16}$, whereas full-target reprojection gives
$1.77\,10^{-27}$. The persistent energy is physical within the discrete
state, not merely an artifact of the subtractive diagnostic.
\begin{figure}[htbp]\centering
\includegraphics[width=.8\linewidth]{night-positive-energy.pdf}
\caption{Positive common-quadrature energy exposes the residual
component in \eqref{eq:increment}. Instrumented repetitions are excluded
from performance timings.}
\end{figure}

\paragraph{Cost and nonuniform deformation.}
Cooperative CUDA quadrature reduces cross-mass assembly from roughly
$0.922$ to $0.122$ seconds on the tested $n_{FE}=256$, $n_W=128$
configuration. A complete adaptive deform-and-return calculation with
255 support changes takes about $55$ rather than $255$ seconds.
The final errors vary slightly across repetitions because atomics and
nearly singular small Grams feed the adaptive geometry. The report
retains all repetitions and their error ranges.

The graph variant is not uniformly superior: on the long deform-and-return
case it gives error $0.0279$ in about $99$ seconds with inverse degree 16,
against about $0.0233$ in $55$ seconds for physical Galerkin memory.
On a prescribed identical sequence of 63 supports, graph degrees 16
and 20 differ by at most $4.1\,10^{-11}$ in component FE norm.
This separates inverse accuracy from adaptive mesh feedback.

\paragraph{What remains necessary for an article.}
The principal unresolved issues are a consistency estimate for general
FE velocities, a coupled PDE example where persistent memory is
indispensable, robust adaptation for multiple separated structures,
and comparison with closer GPU remapping and multiscale baselines.
The present construction and evidence support further study; they do
not establish a revolutionary or publication-ready result.

\section*{Reproducibility}
Branch \texttt{SLFE} of \url{https://github.com/keckj/rodin} contains the
sources under \texttt{experiments/slfe}. The \texttt{results\_night2}
records contain parameters, provenance and per-run diagnostics.
The main scripts are \texttt{remap\_consistency\_study.py},
\texttt{reaction\_moments\_study.py},
\texttt{spectral\_reaction\_study.py},
\texttt{increment\_conditioning\_study.py} and
\texttt{long\_return\_benchmark.py}.
The illustrated report distinguishes physical PDE accuracy, algebraic
oracles, exploratory timings and unsuccessful variants.

\begin{thebibliography}{9}
\bibitem{ouardghi} A. Ouardghi, M. El-Amrani and M. Seaid (2022).
An adaptive enriched semi-Lagrangian finite element method for coupled flow-transport problems.
\href{https://doi.org/10.1016/j.compfluid.2022.105474}{Computers \& Fluids, 105474}.
\bibitem{codina} R. Codina, J. Principe, O. Guasch and S. Badia (2007).
Time dependent subscales in the stabilized finite element approximation of incompressible
flow problems. \href{https://doi.org/10.1016/j.cma.2007.01.002}{CMAME}.
\bibitem{simon} K. Simon and J. Behrens (2021).
Semi-Lagrangian Subgrid Reconstruction for Advection-Dominant Multiscale Problems with Rough Data.
\href{https://doi.org/10.1007/s10915-021-01451-w}{Journal of Scientific Computing}.
\bibitem{adams} M. F. Adams, D. S. Finn, M. G. Knepley and J. V. Pusztay.
A projection method for particle resampling.
\href{https://arxiv.org/abs/2501.13681}{arXiv:2501.13681}.
\bibitem{keith} B. Keith et al. (2026).
Remap for ALE Hydrodynamics via Direct Interpolation and Constraint Optimization.
\href{https://arxiv.org/abs/2611.05327}{arXiv:2611.05327}.
\bibitem{gaul} A. Gaul, M. Gutknecht, J. Liesen and R. Nabben (2013).
A framework for deflated and augmented Krylov subspace methods.
\href{https://people.math.ethz.ch/~mhg/pub/defl-GGLN1-SIMAX.pdf}{SIAM Journal on Matrix Analysis and Applications}.
\bibitem{cayley} E. Mac\'ias-Virg\'os, M. J. Pereira-S\'aez and D. Tanr\'e (2018).
Cayley transform on Stiefel manifolds.
\href{https://doi.org/10.1016/j.geomphys.2017.08.011}{Journal of Geometry and Physics 123, 53--60}.
\end{thebibliography}
\end{document}
